
\subsection{Rescaling Modulo of Congruence Relations}
\label{subsec:modulo-rescaling}

Let $q,q'$ be positive integer moduli, $r=q'/q$, and fix the translation-equivariant nearest-integer rule $R(x)=\lfloor x+1/2\rfloor$. Thus $R(x+h)=R(x)+h$ for integer $h$. Write $\delta_x=R(rx)-rx$, so $|\delta_x|\le1/2$.

\begin{table}[h!]
\centering
\begin{tabular}{|l|l|}
\hline
Original relation & Scaled relation modulo $q'$ \\
\hline
$a\equiv b\pmod q$ & $R(ra)\equiv R(rb)\pmod{q'}$ \\
$a+c\equiv b+d\pmod q$ & $R(ra)+R(rc)\equiv R(rb)+R(rd)+\eta\pmod{q'}$ \\
$ac\equiv bd\pmod q$ & $R(rac)\equiv R(rbd)\pmod{q'}$ \\
\hline
\end{tabular}
\caption{Scaling congruences, where $\eta=\delta_a+\delta_c-\delta_b-\delta_d$ and $|\eta|\le2$.}
\label{tab:rescaling}
\end{table}

\begin{proof}
If $a=b+kq$, then $ra=rb+kq'$ and $R(ra)=R(rb)+kq'$. Thus rounding an entire congruence preserves it exactly, whether or not $ra$ is an integer. The same argument applies to the entire products $ac$ and $bd$.

For separately rounded summands, $a+c=b+d+kq$ gives
\[
R(ra)+R(rc)-R(rb)-R(rd)=kq'+\delta_a+\delta_c-\delta_b-\delta_d.
\]
The last four terms form the integer $\eta$, with $|\eta|\le2$. If they cancel (in particular, if all scaled summands are integers), the relation is exact. Otherwise this explicit error must be retained.
\end{proof}

These statements are forward implications, not general equivalences. For instance, with $q=8$ and $q'=2$, both $a=0$ and $b=1$ round to zero, although $a\not\equiv b\pmod8$. Rounding loses information. Also, separately scaling both factors of a product changes its scale to $r^2$; it cannot replace scaling the entire product once by $r$.

For a fixed integer multiplier $s$, the useful identity is $R(ra)s=ras+\delta_a s$, with error at most $|s|/2$. In an LWE phase $b-\sum_i a_i s_i$, separately rounding $b$ and the $a_i$ therefore introduces $\delta_b-\sum_i\delta_{a_i}s_i$. This error, rather than an undefined approximate-congruence symbol, is what later modulus-switching correctness must bound.

\subsubsection{Example}
\label{subsec:modulo-rescaling-ex}

Suppose we have the following congruence relation:

$b \equiv a\cdot s + \Delta \cdot m + e \bmod q$, \text{ } where: $q = 30$, \text{ } $s = 5$, \text{ } $a = 10$, \text{ } $\Delta = 10$, \text{ } $m = 1$, \text{ } $e = 10$, \text{ } $b = 40$

$ $

First, we can test if the above congruence relation is true by plugging in the given example values as follows: 

$b \equiv a\cdot s + \Delta \cdot m + e \bmod 30$

$40 \equiv 10 \cdot 5 + 10 \cdot 1 + 10 \bmod 30$

$40 \equiv 70 \bmod 30$ 

$ $

This congruence relation is true. 

$ $

Now, suppose we want to rescale the modulo from $30 \rightarrow 3$. Then, based on the rescaling principles described in \autoref{tab:rescaling}, we compute the rescaled values as follows: 

$q'= 3$, \text{ } $s = 5$, \text{ } $m = 1$

$\hat{a} = \Bigg\lceil a\cdot\dfrac{3}{30} \Bigg\rfloor = \Bigg\lceil 10\cdot\dfrac{3}{30} \Bigg\rfloor = 1$

$\hat{\Delta} = \Bigg\lceil \Delta\cdot\dfrac{3}{30} \Bigg\rfloor = \Bigg\lceil 10\cdot\dfrac{3}{30} \Bigg\rfloor = 1$



$\hat{e} = \Bigg\lceil e\cdot\dfrac{3}{30} \Bigg\rfloor = \Bigg\lceil 10\cdot\dfrac{3}{30} \Bigg\rfloor = 1$

$\hat{b} = \Bigg\lceil b\cdot\dfrac{3}{30} \Bigg\rfloor = \Bigg\lceil 40\cdot\dfrac{3}{30} \Bigg\rfloor = 4$

$ $

The rescaled congruence relation from modulo $30 \rightarrow 3$ is derived as follows:

$\Bigg\lceil b\dfrac{3}{30} \Bigg\rfloor \equiv \Bigg\lceil s\cdot a \dfrac{3}{30} \Bigg\rfloor + \Bigg\lceil m \cdot \Delta \dfrac{3}{30} \Bigg\rfloor + \Bigg\lceil e \dfrac{3}{30} \Bigg\rfloor \bmod 3$

$\hat{b} \equiv \hat{a} \cdot s + \hat{\Delta} \cdot m + \hat{e}  \bmod 3$
 \text{ } (an exact congruence relation, as all rescaled values have no decimals)

$4 \equiv 1 \cdot 5 + 1 \cdot 1 + 1 \bmod 3$

$4 \equiv 7 \bmod 3$

$ $

The example is exact because $ra$, $r\Delta$, and $re$ are integers, so every rounding error is zero. If some scaled terms are not integers, exactness is still possible through cancellation; in general, retain the corresponding errors from \autoref{tab:rescaling}.









